\documentclass[12pt,a4paper]{article}
\usepackage[margin=0.75in]{geometry}
\usepackage[utf8]{vietnam}
\usepackage{amsmath, amssymb}
\usepackage{mathtools, bm, cancel, mathrsfs, esint, bbm}
\usepackage[most]{tcolorbox}
\usepackage{xcolor}

\definecolor{myteal}{RGB}{0,75,75}
\definecolor{mybg}{RGB}{242,247,247}
\definecolor{myborder}{RGB}{0,75,75}

\definecolor{titlebg}{RGB}{0,75,75}
\definecolor{titleborder}{RGB}{0,75,75}
\definecolor{titletext}{RGB}{255,255,255}   % màu chữ mong muốn

% --- ĐỊNH NGHĨA KHUNG ---
\tcbset{
    framecauhoi/.style={
        enhanced,
        colback=mybg,
        colframe=myborder,
        boxrule=0.8pt,
        arc=3mm,
        drop shadow={black!40!white},
        left=10pt, right=10pt, top=12pt, bottom=10pt,
        attach boxed title to top left={xshift=10pt, yshift=-10pt},
        coltitle=titletext,
        fonttitle=\bfseries,
        boxed title style={
            colback=titlebg,
            colframe=titleborder,
            boxrule=1pt,
            arc=2mm,
            top=2pt, bottom=2pt, left=6pt, right=6pt,
            drop shadow={black!40!white},
        }
    }
}

% --- HÀM CÂU HỎI ---
\newcommand{\probox}[3][Enter somthing please]{%
    \begin{tcolorbox}[framecauhoi, title={#1}]
        #2
    \end{tcolorbox}
    
    \noindent{\bfseries\color{myteal} LỜI GIẢI}
    
    \vspace{-4pt}
    #3
    \vspace{1.5em}
}

\setlength{\parindent}{0pt}
\setlength{\parskip}{1em}  

\begin{document}

\probox[Bài toán]{Tìm tất cả các hàm số $f:\mathbb{R}^+ \rightarrow \mathbb{R}^+$ thỏa mãn
$$\sqrt{\frac{x^2 + f(y)^2}{2}}\ge \frac{f(x)+y}{2}\ge \sqrt{xf(y)}, \ \forall x, y\in\mathbb{R}^+$$
}{
Trước hết $g(t) := f(t) - t$ \\
Từ đề bài ta có:
$$
2\left(x^2+f(y)^2\right) \ge \left(f(x)+y\right)^2 \ge 4xf(y), \ \forall x,y\in\mathbb{R}^+ \eqno{(\star)}
$$

\textbf{\textcolor{myteal}{Khẳng định 1:}} Với mọi hằng số $c \ge 0$, hàm số $f(x) = x + c$ thỏa mãn $(\star)$ 

Cố định một $c$ như vậy, đặt $a = x$, $b = f(y) = y+c$. Khi đó $(\star)$ trở thành
$$2(a^2+b^2) \ge (a+b)^2 \ge 4ab$$
Bất đẳng thức trên quy về $(a-b)^2 \ge 0$ (luôn đúng). \\
Do đó $f(x)=x+c$ là một nghiệm với mọi $c\ge 0$; điều kiện $c\ge0$ chính là điều kiện cần để \\
$f(x)=x+c>0$ đúng với mọi $x\in\mathbb{R}^+$.

\textbf{\textcolor{myteal}{Khẳng định 2:}} $f(f(y)) = 2f(y) - y, \ \forall y \in \mathbb{R}^+$

Thay $x = f(y)$ vào $(\star)$ ta có:\\
$$
\begin{aligned}
&2f(y)^2 \ge (f(f(y))+y)^2 \ge 4f(y)^2 \\
\implies &f(y) \ge \frac{f(f(y))+y}{2} \ge f(y)\\
\implies &f(f(y)) = 2f(y) - y, \ \forall y \in \mathbb{R}^+
\end{aligned}
$$



\textbf{\textcolor{myteal}{Khẳng định 3:}} $f(y) \ge y, \ \forall y \in \mathbb{R}^+$

Cố định $y \in \mathbb{R}^+$ và định nghĩa dãy $(y_n)$ được xác định bởi $y_0 = y,\ y_{n+1} = f(y_n)$ với $n \ge 0$; vì $f$ ánh xạ từ $\mathbb{R}^+$ vào chính nó nên mỗi $y_n$ là một số thực dương.\\
Áp dụng Khẳng định 2 thu được:
$$ 
\begin{aligned}
&y_{n+1} = f(y_n) = f(f(y_{n-1})) = 2f(y_{n-1}) - y_{n-1} = 2y_n - y_{n-1}\\
\implies &y_n = y_0 + n(y_1 - y_0) = y + ng(y), \ \forall n\in\mathbb{N}
\end{aligned}
$$
Nếu $g(y) < 0$, thì $y_n \to -\infty$ khi $n\to +\infty$, mâu thuẫn với $y_n \in \mathbb{R}^+, \ \forall n\in\mathbb{N}$. Do đó $g(y) \ge 0$, tức là $f(y) \ge y$.

\textbf{\textcolor{myteal}{Khẳng định 4:}} $\lvert g(z) - g(y) \rvert \le \dfrac{(z-y)^2}{4\min\left\{ f(y),f(z)\right\} }, \ \forall y,z \in \mathbb{R}^+$

Thay $x = f(z)$ vào $(\star)$, và sử dụng Khẳng định 2:
$$2\left(f(z)^2+f(y)^2\right) \ge \left(2f(z)-z+y\right)^2 \ge 4f(z)f(y)$$
Khai triển hai vế bất đẳng thức bên phải:
$$\left(z+2g(z)+y\right)^2 \ge 4(z+g(z))(y+g(y))$$
Khai triển cả hai vế và áp dụng Khanqgr định 3 cho ta
$$(z-y)^2 + 4f(z)\big(g(z)-g(y)\big) \ge 0$$
tức là,
$$g(y) - g(z) \le \dfrac{(z-y)^2}{4f(z)}$$
Vì bất đẳng thức này đúng với mọi cặp $(y,z) \in \mathbb{R}^+\times \mathbb{R}^+$, đổi vai trò $y \leftrightarrow z$ cho ta chặn còn lại
$$g(z) - g(y) \le \frac{(z-y)^2}{4f(y)}$$
\\
Như vậy khẳng định được chứng minh

\textbf{\textcolor{myteal}{Khẳng định 5:}} $g$ là hàm hằng trên $\mathbb{R}^+$

Cố định $z,y \in \mathbb{R}^+$ với $z < y$ khi đó $z = \min\left\{z,y\right\} > 0$. Với $n \in \mathbb{Z}^+$, đặt $z_k = z + \frac{k}{n}(y-z)$ với $k = 0,1,\dots,n$, khi đó $z_0 = z$, $z_n = y$, và $z_k \ge z$ với mọi $k$. Theo Khẳng định 3, $f(z_k) \ge z_k \ge z$, do đó áp dụng Khẳng định 4 cho các điểm liên tiếp cho ta
$$\lvert g(z_{k-1}) - g(z_{k}) \rvert \le \frac{(z_{k}-z_{k-1})^2}{4m} = \frac{(y-z)^2}{4zn^2}$$
Theo đó với $k=0,\dots,n-1$ thì:
$$0\le \lvert g(z)-g(y)\rvert = \sum_{k=0}^{n-1} \lvert g(z_{k-1})-g(z_{k})\rvert \le n \cdot \frac{(y-z)^2}{4zn^2} = \frac{(y-z)^2}{4zn}$$
Từ đây cho $n \to +\infty$ suy ra $g(z) = g(y)$. Vì $z,y$ bất kỳ, $g$ là hàm hằng trên $\mathbb{R}^+$.

\textbf{\textcolor{myteal}{Kết luận:}} Theo Khẳng định 5, $g(x) = f(x) - x \equiv c$ với một hằng số $c$ nào đó, và theo Khẳng định 3, $c = g(x) \ge 0$ với mọi $x$. Như vậy mọi nghiệm đều có dạng $f(x) = x+c$ với $c \ge 0$, và theo Khẳng định 1, mọi hàm như vậy thực sự là nghiệm. Do đó $$f(x) = x + c, \ \text{với hằng số tùy ý } c \ge 0$$ là tập hợp tất cả các nghiệm.
\hfill $\square$
\end{document}
}

\end{document}
